\section{The smallest coordinate and the strip exponent}\label{sec:hard-edge}

The positive integral reduces the growth of the arch moment to the coordinates nearest zero. We prove the required bounds by comparing with an explicitly normalized Bures--Laguerre density. The comparison is at the level of positive measures, so it also controls the possible sign changes in the product appearing inside the integral.


The integral formula immediately gives
\begin{equation}\label{eq:01-inverse-comparison}
 \E y_1^{-1/4}\lesssim m_N
 \lesssim\E\left(\sum_{i=1}^Ny_i^{-1}\right)^{1/4}.
\end{equation}
To verify both estimates pointwise, put
$R(x)=\prod_i(y_i-x)/(y_i+x)$. Since each factor lies in $[-1,1]$,
$1-R(x)\ge0$, and telescoping the product gives
\[
 1-R(x)\le2\min\left\{1,x\sum_i y_i^{-1}\right\}.
\]
For $y_1/3<x<y_1/2$, all factors are nonnegative and the first is at
most $1/2$, so $1-R(x)\ge1/2$. Finally,
$\omega(x)\lesssim x^{-1/4}$ on $(0,1)$ and
$\omega(x)\asymp x^{-1/4}$ on $(0,1/2)$. Integrating the lower bound
on the indicated interval and splitting the upper integral at
$x=(\sum_i y_i^{-1})^{-1}$ proves
\eqref{eq:01-inverse-comparison}.

The reference inverse-moment estimates below will contribute a factor
$N^{1/2}$. Scaling those coordinates by a factor of order $N^{-1}$
will contribute a further $N^{1/4}$, producing the required $N^{3/4}$
for the arch moment. We now prove both the reference estimates and the
comparison that permits this scaling.

For comparison we use the generalized Bures--Laguerre family
\cite{ForresterKieburg2016}, in ordered coordinates and with parameter
$a>0$, having density
\begin{equation}\label{eq:01-Laguerre-law}
 \frac1{Z_n(a)}\prod_{i=1}^n r_i^{a-1}e^{-r_i}
          \prod_{i<j}\frac{(r_j-r_i)^2}{r_i+r_j},
 \qquad 0<r_1<\cdots<r_n<\infty.
\end{equation}
Expectations and probabilities for this law will carry subscripts
$n,a$ when clarification is needed.
The exponent of the one-variable weight is $a-1$ in our convention;
thus our $a$ equals the parameter in \cite[equation~(1.5)]{ForresterKieburg2016}
plus one. Our ordered normalizer is the corresponding unordered integral
divided by $n!$.
The following association statement is the continuous MTP$_2$ principle
of Sarkar, for log-supermodular densities
\cite[Theorem 2.27]{Bogso2014}; see also
\cite[equation (3.5)]{FallatEtAl2017}. We give a direct conditioning
proof for the ordered pair kernel used here.

\begin{lemma}[Association and monotone comparison]\label{lem:association}
Let $I\subset(0,\infty)$ be an interval, let $n\ge1$, and let
$w_1,\ldots,w_n$ be positive smooth functions on $I$. Suppose
\[
 \prod_{i=1}^n w_i(x_i)\prod_{i<j}K(x_i,x_j),\qquad
 K(x,y)=\frac{(y-x)^2}{x+y},\qquad x_1<\cdots<x_n,\quad x_i\in I,
\]
has finite positive normalizing integral. Its probability law is
positively associated: increasing integrable functions have nonnegative
covariance whenever the covariance is defined. The assertion also
holds after truncation to a shorter common interval. Consequently,
a coordinatewise increasing likelihood ratio shifts the coordinates
upwards in stochastic order, and a decreasing one shifts them
downwards.
\end{lemma}

\begin{proof}
The mixed derivative of each pair contribution is
\[
 \frac{\partial^2}{\partial x\partial y}\log K(x,y)
 =\frac2{(y-x)^2}+\frac1{(x+y)^2}>0.
\]
We prove association by induction on the number of coordinates,
initially for bounded increasing functions. Conditional on the
largest coordinate $t$, the remaining coordinates have a law of the
same form on the shorter interval, with modified one-variable
factors. They are associated by induction. If $s<t$, the ratio of
the conditional density at $s$ to that at $t$, on the former support,
is coordinatewise decreasing by the mixed-derivative inequality.
Extended by zero where a coordinate exceeds $s$, it remains
decreasing on the latter support. Association for the conditional
law at $t$ shows that a decreasing likelihood tilt decreases every
increasing expectation. If the ratio is unbounded, first truncate
it and pass to the limit by monotone convergence. Thus the
conditional expectation of any increasing function of all the
coordinates increases with $t$.

The law of total covariance now has two nonnegative terms:
conditional covariance by induction, and covariance of the two
conditional expectations by the one-dimensional inequality
\[
 \operatorname{Cov}(F(T),G(T))
 =\frac12\E[(F(T)-F(T'))(G(T)-G(T'))]\ge0.
\]
Here $T'$ is an independent copy of $T$. This proves association;
truncation and integrable approximation extend the assertion to the
stated class. Finally for an increasing likelihood ratio $L$,
$\E[FL]/\E L\ge\E F$ for increasing $F$, and the reversed inequality
holds for decreasing $L$. This is the claimed stochastic comparison.
\end{proof}

The upper comparison will use $a=1/2$, and the lower comparison will
use $a=1$. To estimate the minimum for $a=1$, we will first need an
inverse-moment bound for $a=2$. The next lemma develops all three
estimates from the same normalizing integral.

\begin{lemma}[Laguerre normalization and coordinate bounds]
\label{lem:laguerre}
For fixed $a>0$ and every integer $n\ge1$,
\begin{equation}\label{eq:01-normalizer}
 Z_n(a)=\prod_{i=0}^{n-1}\Gamma(a+i)
             \prod_{0\le i<j<n}\frac{j-i}{2a+i+j}.
\end{equation}
There exist constants $K_a,C_a$, depending only on $a$, such that
$\Prob_{n,a}(r_n\le K_an)$ is bounded below by a positive constant
independent of $n$, and, with $D_k=ak+k(k-1)/2$,
\begin{equation}\label{eq:01-small-coordinate}
 \Prob_{n,a}(r_k<h)
 \le\min\left\{1,\left(\frac{C_a h n^2}{k^3}\right)^{D_k}\right\},
 \qquad h>0,\quad 1\le k\le n.
\end{equation}
In particular,
\begin{align}
 \E_{n,1/2}\left(\sum_i r_i^{-1}\right)^{1/4}
       &\lesssim n^{1/2},\label{eq:01-reference-upper}\\
 \E_{n,2}\sum_i r_i^{-1}&\lesssim n^2,
       \label{eq:01-reference-inverse}\\
 \E_{n,1}r_1^{-1/4}&\gtrsim n^{1/2}.
       \label{eq:01-reference-lower}
\end{align}
The constants in the last three bounds are absolute.
\end{lemma}

\begin{proof}
\emph{Normalization.}
Write the pair product as a Vandermonde determinant times the Schur
Pfaffian with entries $(s-r)/(s+r)$. The classical identity, in the
form recorded in \cite[equation (1.2)]{IshikawaOkadaTagawaZeng2006}, is
\[
 \operatorname{Pf}\left[\frac{x_j-x_i}{x_j+x_i}\right]
 =\prod_{i<j}\frac{x_j-x_i}{x_j+x_i}.
\]
For even size, clearing denominators leaves an alternating
polynomial of Vandermonde degree, so only its constant factor needs
checking. Pairing two opposite arguments gives that constant
recursively as one. For odd size, add an argument tending to infinity;
the resulting augmenting column consists of ones.

Expand the determinant and the Pfaffian and integrate over the
unordered region, divided by $n!$. This is de Bruijn's
determinant--Pfaffian integration formula~\cite{deBruijn1955}; the
general antisymmetric-kernel form, including its odd-size border,
is stated in \cite[Theorem 6.1]{BaikRains2001Algebraic}. The result is
the Pfaffian of double moments, with the single moments as augmenting
entries in odd size. Absolute integrability follows from
$|(s-r)/(s+r)|\le1$ and the finite gamma moments for $a>0$.
The double moment with indices $i,j$ is
\[
 \int_0^\infty\!\int_0^\infty
 r^{a+i-1}s^{a+j-1}e^{-r-s}\frac{s-r}{s+r}\,dr\,ds
 =\Gamma(a+i)\Gamma(a+j)\frac{j-i}{2a+i+j}.
\]
Use $r+s$ and $r/(r+s)$ to evaluate the integral; the latter beta
integral has mean $(a+i)/(2a+i+j)$. Factoring out the gamma terms and
using the Schur identity again, now at the arguments $a+i$, proves
\eqref{eq:01-normalizer}.

\smallskip\noindent\emph{Coordinate tails and inverse moments.}
For the maximum, condition on its value $s$. Every factor involving
a smaller coordinate is at most $s$, since $K(r,s)\le s$ for
$0<r<s$. Consequently
\[
 \Prob_{n,a}(r_n>Kn)
 \le\frac{Z_{n-1}(a)}{Z_n(a)}
         \int_{Kn}^\infty s^{a+n-2}e^{-s}\,ds.
\]
The explicit normalizer gives
\[
 \frac{Z_{n-1}(a)}{Z_n(a)}
 =\frac1{\Gamma(a+n-1)}
   \prod_{j=1}^{n-1}\frac{2a+2n-2-j}{j}
 \le\frac{e^a(2e)^{n-1}}{\Gamma(a+n-1)}.
\]
For $n\ge2$, the inequality follows by replacing every numerator by
$2(a+n-1)$, using $(n-1)!\ge((n-1)/e)^{n-1}$, and then
$(1+a/(n-1))^{n-1}\le e^a$; $n=1$ is immediate. The exponential-moment bound for a gamma random
variable gives
\[
 \frac1{\Gamma(a+n-1)}\int_{Kn}^\infty s^{a+n-2}e^{-s}\,ds
 \le e^{-Kn/2}2^{a+n-1}.
\]
Choosing $K$ large makes the desired tail uniformly small, and in
fact exponentially small in $n$.

For the small-coordinate estimate, integrate the first $k$
coordinates below $h$. Their mutual pair factors are at most $h$;
a factor between a small coordinate and a remaining one is at most
the latter. Discarding order constraints between these two groups
therefore gives
\begin{equation}\label{eq:01-small-normalizer-ratio}
 \Prob_{n,a}(r_k<h)
 \le\frac{h^{D_k}}{a^k k!}\frac{Z_{n-k}(a+k)}{Z_n(a)},
\end{equation}
with $Z_0=1$. From \eqref{eq:01-normalizer},
\[
 \frac{Z_{n-k}(a+k)}{Z_n(a)}
 =\prod_{i=0}^{k-1}\Gamma(a+i)^{-1}
  \prod_{\substack{0\le i<j<n\\i<k}}
       \left(1+\frac{2a+2i}{j-i}\right).
\]
For gaps $j-i\le k$, use $2a+2i\le C_a k$. For each $i<k$,
\[
 \sum_{j=1}^k\log(1+C_a k/j)
 \le k\log(1+C_a)+k\log k-\log(k!)=O_a(k).
\]
Thus the short gaps contribute $O_a(k^2)=O_a(D_k)$. For larger
gaps, $\log(1+x)\le x$ gives a contribution at most
\[
 \sum_{i<k}(2a+2i)\sum_{j=k+1}^{n-1-i}\frac1j
 \le 2D_k\log(n/k).
\]
This includes $k=n$, when all the latter sums are empty.
Furthermore
\begin{equation}\label{eq:01-gamma-product}
 \prod_{i=0}^{k-1}\Gamma(a+i)\ge k^{D_k}e^{-O_a(D_k)}.
\end{equation}
Here is a direct estimate. For $q=a+i$, integration on $[q,q+1]$
gives
$\log\Gamma(q)\ge(q-1)\log q-q-O_a(1)$.
Summing, separating the factor $\log k$, and using
$\sum_{i<k}(a+i)\log((a+i)/k)\ge-k^2/e$ gives
$D_k\log k-O_a(D_k)$; the missing $k\log k$ is absorbed in
$O_a(D_k)$. Adjust the constant for bounded $k$.
Combining these bounds with
\eqref{eq:01-small-normalizer-ratio}, and absorbing $a^{-k}/k!$
into $C_a^{D_k}$, proves \eqref{eq:01-small-coordinate}.

For any $0<p<D_k$, integration of that tail estimate gives
\begin{equation}\label{eq:01-inverse-tail}
 \E_{n,a}r_k^{-p}
 \le\frac{D_k}{D_k-p}
       \left(\frac{C_a n^2}{k^3}\right)^p.
\end{equation}
At $a=1/2$, apply this with $(k,p)=(1,1/4)$ and then with $p=1$,
$k\ge2$. It gives
$\E r_1^{-1/4}\lesssim n^{1/2}$ and
$\E\sum_{k\ge2}r_k^{-1}\lesssim n^2$.
Subadditivity of the fourth root and Jensen's inequality yield
\eqref{eq:01-reference-upper}. At $a=2$, use $p=1$ for every $k$
to obtain \eqref{eq:01-reference-inverse}.

\smallskip\noindent\emph{The minimum at $a=1$.}
The density of the smallest coordinate at $s>0$ is
\[
 e^{-s}\frac{Z_{n-1}(2)}{Z_n(1)}
 \E_{n-1,2}\left[
    \mathbf1_{\{r_i>s\ \forall i\}}
    \prod_i\frac{(1-s/r_i)^2}{1+s/r_i}\right].
\]
The normalizer ratio is exactly $n(n+1)/2$.
Set $q_s(r)=\mathbf1_{\{r>s\}}(1-s/r)^2/(1+s/r)$.
For all $r>0$, including $r\le s$, one has
$0\le q_s(r)\le1$ and $1-q_s(r)\le3s/r$. Consequently
\[
 \prod_iq_s(r_i)\ge1-3s\sum_i r_i^{-1}.
\]
By \eqref{eq:01-reference-inverse}, the expectation is at least
$1/2$ whenever $0<s<\varepsilon n^{-2}$ for a sufficiently small
fixed $\varepsilon>0$. The displayed density is then bounded below
by a positive multiple of $n^2$. An interval with endpoints fixed
positive multiples of $n^{-2}$ has probability bounded below;
on it $r_1^{-1/4}\gtrsim n^{1/2}$. This proves
\eqref{eq:01-reference-lower}, with $n=1$ also covered directly.
\end{proof}

\begin{proof}[Completion of Theorem~\ref{thm:strip}]
We compare \eqref{eq:01-positive-law} with the reference laws to
estimate the two sides of \eqref{eq:01-inverse-comparison}.
For the upper estimate, first condition the target law on
$y_N<1/2$. This decreasing event shifts all coordinates downwards,
by Lemma~\ref{lem:association}, and therefore increases the inverse
statistic. On this shorter interval compare it with
$y_i=r_i/(bN)$ under \eqref{eq:01-Laguerre-law} with $a=1/2$,
conditioned on $r_N<bN/2$. Up to a constant, the target-to-reference
likelihood ratio is the product of
\[
 \omega(y)y^{1/2}e^{bNy}.
\]
Direct differentiation gives
\[
 \frac{d}{dy}\log(\omega(y)y^{1/2})
 =\frac{1-2y-\sqrt{2y/(1+y)}}{4y(1-y)}\ge-4,
 \qquad 0<y\le\frac12.
\]
Indeed the numerator is positive for $y\le1/8$; on the remaining
interval it is at least $-1$ and the denominator is at least $1/4$.
Thus the likelihood factors are increasing for every $N\ge1$ once
$b\ge4$. Association shows that the conditioned reference
law has the larger inverse statistic. Choosing also $b/2\ge K_{1/2}$
makes its conditioning probability bounded below uniformly in $N$,
by Lemma~\ref{lem:laguerre}. Hence
\[
 \E\left(\sum_i y_i^{-1}\right)^{1/4}
 \lesssim (bN)^{1/4}
          \E_{N,1/2}\left(\sum_i r_i^{-1}\right)^{1/4}
 \lesssim N^{3/4}.
\]

For the lower estimate, set $y_i=2z_i/(1+z_i)$ in the target law,
so $0<z_1<\cdots<z_N<1$. The transformed pair kernel is
\[
 K\left(\frac{2z_i}{1+z_i},\frac{2z_j}{1+z_j}\right)
 =\frac{2K(z_i,z_j)}
 {(1+z_i)(1+z_j)(1+2z_iz_j/(z_i+z_j))}.
\]
Together with the Jacobian this gives one-variable factors
$\omega(2z/(1+z))(1+z)^{-(N+1)}$ and additional pair factors
$(1+2z_iz_j/(z_i+z_j))^{-1}$.
Compare with $z_i=r_i/(bN)$ under the reference law with $a=1$ and
fixed $0<b<1/2$, conditioned on $r_N<bN$. The ratio has individual
factors
\[
 \omega\left(\frac{2z}{1+z}\right)(1+z)^{-(N+1)}e^{bNz}
\]
and the additional pair factors just displayed. Each is decreasing
in every coordinate: $\omega$ is decreasing, and
$-(N+1)/(1+z)+bN<0$ for $0<z<1$. Association therefore puts the
target $z$ coordinates below those of this conditioned reference,
which are themselves below the unconditioned reference coordinates.
Since $y_1\le2z_1$, it follows that
\[
 \E y_1^{-1/4}
 \ge 2^{-1/4}(bN)^{1/4}\E_{N,1}r_1^{-1/4}
 \gtrsim N^{3/4}.
\]
Equation \eqref{eq:01-inverse-comparison} now proves
$m_N\asymp N^{3/4}$.

It remains to extract the bridge estimate from
Lemma~\ref{lem:moment-increments}. By monotonicity, summing the lower
increment bound over $1\le j<N$ gives
$(N-1)\mathcal B_N\lesssim m_N$, so
$\mathcal B_N\lesssim N^{-1/4}$, with bounded $N$ absorbed in the
constant. Conversely choose a fixed integer $L$ so large that the
lower bound for $m_{LN}$ exceeds twice the upper bound for $m_N$.
Then
\[
 N^{3/4}\lesssim m_{LN}-m_N
 \lesssim\sum_{j=N}^{LN-1}\mathcal B_j
 \le(L-1)N\mathcal B_N.
\]
This proves $\mathcal B_N\gtrsim N^{-1/4}$ and completes all assertions
of the theorem.
\end{proof}
